% Generated draft: requires mathematical and source review.
\documentclass[11pt]{article}
\usepackage[margin=25mm]{geometry}
\usepackage{fontspec}
\setmainfont{Latin Modern Roman}
\usepackage{amsmath,amssymb,mathtools,unicode-math}
\setmathfont{Latin Modern Math}
\usepackage{enumitem,needspace,xurl,hyperref}
\setlength{\parindent}{0pt}
\setlength{\parskip}{5pt}
\setlength{\emergencystretch}{3em}
\newcommand{\N}{\mathbb{N}}
\newcommand{\Z}{\mathbb{Z}}
\newcommand{\Q}{\mathbb{Q}}
\newcommand{\R}{\mathbb{R}}
\newcommand{\C}{\mathbb{C}}
\newcommand{\F}{\mathbb{F}}
\newcommand{\A}{\mathbb{A}}
\newcommand{\PP}{\mathbb{P}}
\newcommand{\E}{\mathbb{E}}
\newcommand{\eps}{\varepsilon}
\newcommand{\dd}{\,\mathrm{d}}
\begin{document}
\section*{UnsolvedMath ID: 800914}
\textbf{Source problem code:} catalogue entry
\subsection{Definitions and mathematical statement}
The supplied record does not specify the hypotheses on \(X\), the domain of \(\beta\), the genus range, or the coefficient rings. The following gives a conventional setting in which its displayed questions are defined; this setting requires confirmation against the source.

Let \(X\) be a smooth projective variety over \(\C\), let \(g\geq 2\), and let \(\beta\in H_2(X,\Z)\) be a curve class. Let \(\overline{M}_g(X,\beta)\) be the moduli stack of connected, unmarked, genus-\(g\) stable maps to \(X\) representing \(\beta\), with virtual fundamental class \([\overline{M}_g(X,\beta)]^{\operatorname{vir}}\). Let
\[
\pi:\overline{M}_g(X,\beta)\longrightarrow\overline{M}_g
\]
forget the map and stabilize the domain, where \(\overline{M}_g\) is the moduli stack of stable curves of genus \(g\). Use rational Chow groups and rational cohomology, identifying the homological pushforward of the virtual class with its corresponding cohomology class on \(\overline{M}_g\).

For stable pairs \((h,n)\), the tautological Chow rings \(R^*(\overline{M}_{h,n})\subseteq A^*(\overline{M}_{h,n};\Q)\) form the smallest system of unital \(\Q\)-subalgebras closed under pushforward by the standard marking-forgetting and boundary-gluing morphisms, including gluing from products. Write \(RH^*(\overline{M}_{h,n})\) for their images under the cycle-class map.

\textbf{Q1 (the primary source item).} Is it true, for every \((X,g,\beta)\) in the intended admissible range, that
\[
\operatorname{cl}\!\left(\pi_*[\overline{M}_g(X,\beta)]^{\operatorname{vir}}\right)
\in RH^*(\overline{M}_g)?
\]
The supplied excerpt does not determine whether the conventional range stated above is exactly the intended range.

\textbf{Q2 (also included in the supplied record).} Determine when
\[
\pi_*[\overline{M}_g(X,\beta)]^{\operatorname{vir}}
\in R^*(\overline{M}_g)
\subseteq A^*(\overline{M}_g;\Q).
\]

The source expresses an expectation that Q1 has an affirmative answer. It states that Q1 is affirmative when \(X\) is a curve, citing the unresolved reference label \texttt{[fabp]}, and that Q2 has counterexamples even when \(X\) is a curve. These are supplied source assertions, not independently verified conclusions here. The source further presents as a \emph{wild speculation}, motivated by the Bloch--Beilinson conjecture, that the Chow membership in Q2 holds whenever \(X\) is defined over \(\overline{\Q}\). In the conventional complex setting, this means that \(X\) admits a model over \(\overline{\Q}\). The excerpt supplies no additional arithmetic condition on \(\beta\).
\subsection{Short English statement}
Ask whether pushing the virtual class of the stable-map space to the moduli space of stable curves always gives a tautological cohomology class, and determine when it gives a tautological Chow class. The source expects an affirmative cohomological answer and speculates about a Chow answer for targets defined over \(\overline{\Q}\). The ambient hypotheses and the cited curve results require source review.
\subsection{Sources}
\begin{itemize}
\item Local UnsolvedMath registry, permanent ID \(800914\), exact source problem code \texttt{catalogue entry}; supplied title: \emph{AGNES Open Problems in Algebraic Geometry --- Q1}. The registry attributes the problem to Pandharipande, proposed in 2010, and labels it \emph{open}; this label is supplied metadata, not an independently verified current assessment.
\item Supplied source-list attribution: Pandharipande, \emph{Open Problems in Enumerative Algebraic Geometry (AGNES)} (2010), source item Q1 (\#1), \url{https://arxiv.org/abs/1004.3259}. The supplied provenance records access on 2026-10-16 and extraction by \texttt{arxiv\_latex}; no independent source visit or verification is asserted here. The extracted passage also contains Q2.
\item The excerpt cites \texttt{[fabp]} for the affirmative cohomological result for curve targets, but supplies no bibliographic expansion. It supplies no specific reference for the asserted Chow counterexamples. The Bloch--Beilinson conjecture is named only as motivation for the arithmetic speculation.
\end{itemize}
\end{document}
